%! Observational Studies — Unit 1, Lesson 1.5 — Lesson Plan
% PILOT SHAPE (2026-09-06): modeled on AP Statistics lesson 1.4 minus its AP Practice page.
% GRADUAL-RELEASE plan (warm-up / guided notes and practice / spoken debrief / close and START
% THE HOMEWORK). The guided notes carry the release ROW BY ROW in one Main Ideas / Notes table
% (2026-09-12 shape): I do / we do in rows 1-4 (problems 1-14) -> we do the Guided Practice row
% (problems 15-18). THERE IS NO INDEPENDENT PRACTICE SET IN THE NOTES — the homework is the
% individual practice, and the period ends with students starting it. Teacher-facing. Breakable
% boxes are `skillbox` with a \boxguard ahead; the two tabularx boxes are `fixedskillbox`.
% Regenerated 2026-09-14 from the PRE-EFFL LEGACY shape: activity/ + exit_ticket/ removed, the
% activity's Riverbend Study Center became the Guided Practice row, the exit ticket's Bayside
% item became the debrief's whole-class cold check, and the tiers were cut, not concatenated.
\documentclass[10pt]{article}
\usepackage{saar-article}
\usepackage{saar-boxes}
\usepackage{graphicx}   % -article does NOT load graphicx
\graphicspath{{images/}}
\newcommand{\TallMath}[1]{$\displaystyle #1\rule[-1.4em]{0pt}{3.2em}$}
\newcommand{\UnitNumberName}{Unit 1: Asking Questions with Data: Study Design \quad}
\newcommand{\LessonNumberName}{Lesson 1.5: Observational Studies}

\begin{document}
\begin{center}
    {\Large\bfseries \CourseName \\
    {\small\bfseries \UnitNumberName \LessonNumberName}}
\end{center}

\begin{tcolorbox}[colback=frost, colframe=cerulean, boxrule=0.9pt, arc=2mm,
    left=2mm, right=2mm, top=1.5mm, bottom=1.5mm]
    \textbf{Primary Objective:} Students will explain what makes a study \emph{observational}
    --- nobody was \emph{assigned} to a group --- plan one with the statistical cycle, compute
    and compare rates for two groups that already exist, name the explanatory and response
    variables, state the result as an \textbf{association}, and justify which conclusions the
    study does and does not support, naming a reversed arrow or a \textbf{lurking variable}
    when it does not.
    \\[2pt]
    \textbf{Standards:} PS.DC.2d \quad$\cdot$\quad PS.DC.1a--b carried forward in the cycle and
    the conclusion items \quad$\cdot$\quad previews PS.DC.3a--b (Lesson 1.6)
    \\[2pt]
    \textbf{Lesson model:} \emph{Gradual release --- I do, we do, you do --- all of it inside
    the guided notes.} There is no group activity and no exit ticket. The notes name the
    vocabulary as they teach it and deliver the instruction \textbf{row by row in one Main
    Ideas / Notes table} (four rows, problems 1--14: you read the definition, mark up the
    display, and work the first problem of each grid; the class works the rest); the Guided
    Practice row (problems 15--18) is worked together with students holding the pen;
    \textbf{the homework is the individual practice}, and students start it in class while you
    circulate. The whole-class debrief is \emph{spoken}.
    \\[2pt]
    \textbf{Note on placement:} every scenario today is a random sample with a $100\%$
    response rate. Bias was Lesson~1.4's job and it is deliberately off the table, so the only
    thing left to explain the gap is that \emph{the groups built themselves}. Do not let
    ``the study was biased'' pass as an answer.
\end{tcolorbox}

\renewcommand{\arraystretch}{1.4}
\boxguard
\begin{skillbox}[Learning Targets \& Key Understandings]{goldbox}
\begin{tabularx}{\linewidth}{@{}X|X@{}}
\textbf{Learning targets (student language)} & \textbf{Key understandings (the ``why'')} \\[2pt]
\hline
\vspace{-6pt}
\begin{itemize}[leftmargin=1.1em, nosep, after=\vspace{2pt}]
  \item I can tell an \textbf{observational study} from an experiment by asking \emph{who
        decided the groups}. \hfill \textit{PS.DC.2d}
  \item I can plan an observational study with the \textbf{statistical cycle} --- formulate,
        collect, represent, analyze. \hfill \textit{PS.DC.2d; 1a}
  \item I can compute a rate for each of two existing groups, report the \textbf{gap in
        percentage points}, and name the \textbf{explanatory} and \textbf{response}
        variables. \hfill \textit{PS.DC.2d}
  \item I can state a finding as an \textbf{association} with no causal verb, and explain why
        an association does not show \textbf{causation} --- naming a reversed arrow or a
        \textbf{lurking variable}. \hfill \textit{PS.DC.2d; 1b}
\end{itemize}
&
\vspace{-6pt}
\begin{itemize}[leftmargin=1.1em, nosep, after=\vspace{2pt}]
  \item \emph{A survey is one type of observational study}, so every study in Unit~1 so far
        has been one. The name is new; the thing is not.
  \item \textbf{Target misconception:} that a clean study proves the claim --- \emph{association
        is causation}. The same $56/80$ against $54/120$ table is produced by three live
        stories: the claim, the same arrow \textbf{backwards}, and a \textbf{lurking variable}
        (retirement, $75\%$ against $15\%$) moving both. Ruling any of them out requires
        \emph{assigning} --- Lesson~1.6.
  \item \textbf{An unbiased study is not a conclusive one.} Lesson~1.4 removed bias; a
        perfectly clean sample still cannot say \emph{why} two groups differ.
  \item \textbf{Second misconception:} that a prediction is safer than a claim. \emph{``If the
        evening swimmers switched\ldots''} is a causal claim in disguise (problem~14).
  \item The cycle does not change from study to study. What changes is what you are licensed
        to write at the end of stage four.
\end{itemize}
\end{tabularx}
\end{skillbox}

\boxguard[24]
\begin{skillbox}[{Vocabulary, Concepts, \& Theorems --- taught directly in the guided notes}]{greenbox}
\small
\renewcommand{\arraystretch}{1.35}
\begin{tabularx}{\linewidth}{@{} >{\bfseries}p{3.1cm} X @{}}
Observational study & A study in which you measure or survey a sample without trying to affect
       anyone. The researcher does not assign anybody to a group. \\
Survey & One type of observational study: the members of a sample are asked questions and
       their answers are recorded. \\
Explanatory variable & The variable we think might explain the other one --- the tail of the
       arrow in the claim being tested. \\
Response variable & The variable we measure to see what happens --- the head of the arrow. \\
Association & Two variables move together: knowing which group an individual is in changes
       what you expect for the other variable. \\
Causation & One variable actually produces the change in the other. A far stronger claim than
       association, and out of reach without assigning. \\
Lurking variable & A variable outside the study that affects both variables in it, and can
       create an association all by itself. \\
Reversed arrow & The response variable is what put individuals into their group in the first
       place, so the association runs the other way. \\
\end{tabularx}
\end{skillbox}

% fixedskillbox never splits, so the phase table stays intact on one page.
% THE PHASE TABLE IS A CONTRACT: 5 / 35 / 10 / 10 — the I do is ~20 of the 35 and the Guided
% Practice ~15; the debrief is spoken; the homework start is the second formative read.
\begin{fixedskillbox}[Lesson at a Glance --- \MeetingLength]{frost}
\renewcommand{\arraystretch}{1.25}
\begin{tabularx}{\linewidth}{@{}l c X X@{}}
\textbf{Phase} & \textbf{Min} & \textbf{Students} & \textbf{Teacher} \\
\hline
Warm-Up & 5 & The two Harbor Point group rates and the pooled rate --- silent & Debrief item~3
  aloud; leave $70\%$, $45\%$, $25$ points, and $55\%$ on the board all period \\
Guided Notes \& Practice & 35 & Fill the vocabulary box as each term is named; work rows 1--4
  of the table (problems 1--14) with the pen in their hand; then the Guided Practice row,
  \emph{Study Center} (15--18), together & \textbf{I do / we do, row by row}: read the
  definition, mark up the display, work problems 1, 5, 9, 11; the class works the rest (20)
  $\to$ \textbf{we do} Guided Practice 15--18 (15) \\
Debrief & 10 & Pens down, papers up --- answer aloud, nothing to fill in & Put the three
  stories back on the board, then take the Bayside cold check and the banned-verb list \\
Close \& Assign & 10 & \textbf{Start the homework in class}, alone and silent & Announce the
  homework and its form (packet or Desmos), circulate for the second formative read, preview
  Lesson~1.6 \\
\end{tabularx}
\end{fixedskillbox}

\begin{fixedskillbox}[Warm-Up --- Activate Prior Knowledge (5 min)]{frost}
\begin{tabularx}{\linewidth}{@{}X X@{}}
\begin{minipage}[t]{\linewidth}\vspace{0pt}
    \textbf{The three items and what each seeds}
    \begin{itemize}[leftmargin=1.1em, itemsep=1pt, topsep=2pt]
      \item \textbf{1.} $56$ of $80$ as a percent. \textbf{Seeds:} the morning rate, $70\%$ ---
            read straight off the two-way table at problem~9 in row~3.
      \item \textbf{2.} $54$ of $120$ as a percent. \textbf{Seeds:} the evening rate, $45\%$,
            and \emph{a second, different denominator}. The two wholes are the whole
            difficulty of this lesson.
      \item \textbf{3.} The gap and the pooled rate. \textbf{Seeds:} the $25$-point gap the
            hook argues about, and the $55\%$ problem~10 scales to all $1{,}500$ members ---
            the ``scale a sample up'' skill from Lesson~1.2.
    \end{itemize}
\end{minipage}
&
\begin{minipage}[t]{\linewidth}\vspace{0pt}
    \textbf{Running it}
    \begin{itemize}[leftmargin=1.1em, itemsep=1pt, topsep=2pt]
      \item \textbf{Debrief item~3 aloud} and leave $70\%$, $45\%$, $25$ points, and $55\%$
            on the board \emph{all period}. The hook needs them thirty seconds after the
            warm-up is collected, and row~3 reads them back.
      \item \textbf{Item~3's third row is the tell.} The wrong answer is $(70+45)\div 2 =
            57.5\%$, which averages two unequal groups; the right path is $110\div 200 =
            55\%$. A student who averages will merge the denominators again on homework~3.
      \item Items 1--2 are Algebra~1 percent work the class has done in 1.2 and 1.4 --- say
            so, so nobody thinks they are being retested. Do not let them expand.
    \end{itemize}
\end{minipage}
\end{tabularx}
\end{fixedskillbox}

\boxguard
\begin{skillbox}[Hook --- before anyone sits down]{frost}
The hook is on \textbf{page 1 of the notes}, under the vocabulary box, and on the board:
\textit{``A random sample of $200$ members. Everyone answered. Every count correct. $70\%$ of
morning swimmers get $7$+ hours of sleep; only $45\%$ of evening swimmers do. The director
prints a flyer: `Swim in the morning --- sleep better!'\,''} \textbf{Ask: is the flyer's claim
safe?} Students circle \emph{safe} or \emph{not safe} and write one line of reason in the hook
box \emph{before} anyone speaks; then take the hand vote and write the split on the board.
\textbf{Close every door the class has learned to open first} --- the sample was random, nobody
was left off the list, everyone answered, every count is correct. Most rooms vote \emph{safe},
because Lesson~1.4 taught them that a clean procedure is the fix, and this is the first time a
clean procedure is not enough.
\textbf{Do not settle it.} Say only that the answer is problem~13, in the \emph{Not Causation}
row, and that they will vote again there. The vote is what makes the crux land: students have
to have committed to \emph{safe} before the notes take it away.
\end{skillbox}

\boxguard
\begin{skillbox}[Guided Notes \& Practice --- I do, then we do (35 min)]{frost}
\textbf{This is the whole lesson.} There is no group activity, and there is no independent
practice set on this page: \textbf{the homework is the individual practice}, started in class.
The notes are \textbf{one Main Ideas / Notes table}: four instruction rows (problems 1--14),
then the Guided Practice row (15--18). \textbf{The rhythm in every row is the same}: read the
definition aloud, mark up the display, work the first problem of the grid yourself, then hand
the rest to the room with the pen in their hand. The vocabulary box on page~1 is filled
\emph{as each term is named}, never in advance. There are no fill-in blanks on the page --- every
answer is written in a problem's own space or in a small framed label box under a display.
\begin{multicols}{2}
\textbf{I do / we do, row by row (about 20 min).}

\emph{Row 1, Assign (problems 1--4).} The whole lesson is one word. Read the definition, then
have the class name the two panels --- left is \emph{watch}, right is \emph{assign}, and the
observational one is the left. Work problem~1 yourself; the class does 2--4. \textbf{Problem~3
is the trap}, and it is the only plan on the page where the director \emph{assigns} swim times:
take a hand count before anyone writes, because students sort by \emph{is it a survey} instead
of \emph{who decided the groups}. Problem~4 is the census --- still observational. Four minutes.

\emph{Row 2, The Cycle (problems 5--8).} This is PS.DC.2d proper, and it should feel like
review because it is: Lesson~1.1's four stages, with 1.3's sampling plan in Collect and 1.4's
bias check standing behind it. \textbf{Do not skip it.} Work problem~5 yourself (the frame, the
random $200$, the $100\%$ response); the class does 6--7. \textbf{Problem~8 is the point of the
row}: nothing in the plan tells a member when to swim, and that one absence is the definition,
written as a plan. Five minutes.

\emph{Row 3, Two Variables (problems 9--10).} The arithmetic is already done --- it is the
warm-up --- so keep it quick and spend the time on \emph{explanatory} and \emph{response},
which is new vocabulary and pays forward to scatterplots in Unit~5. Have the class label the
arrow under the table: tail is the explanatory variable, head is the response. Work problem~9
yourself, the class does~10 ($55\%$ scaled to $825$, \emph{an estimate}). Four minutes.

\emph{Row 4, Not Causation (problems 11--14) --- the crux.} Give story~1, then \textbf{ask the
class to invent story~2 --- the arrow backwards --- before you offer it}; somebody reliably says
``maybe they get up early \emph{because} they sleep well.'' Work problem~11 live. Only then
reveal story~3's retirement row and let the class do problem~12: $60/80 = 75\%$ against
$18/120 = 15\%$, a variable that was never in the study at all. \textbf{Then re-take the hook
vote at problem~13 before you say anything.} Land the printed sentence word by word: an
observational study shows \emph{association}; it cannot show \emph{causation}, because nobody
assigned the groups. Problem~14 is the subtle one --- a prediction about switching is a causal
claim wearing different clothes. Seven minutes.

\columnbreak
\textbf{We do (about 15 min).} The Guided Practice row, \emph{Study Center} (problems 15--18)
--- the same moves the homework will ask alone, on a school instead of a pool, and built so the
obvious reading of the data is not merely unsupported but exactly \emph{backwards}.
\textbf{Students hold the pen; you hold the questions.} Problem~15 goes on them cold (who
decided the groups, and the two variables). Problem~16 is the arithmetic, and the surprise is
the direction: $39/60 = 65\%$ against $72/90 = 80\%$, a $15$-point gap with the \emph{non}-users
higher. \textbf{Problem~17 is the crux transferred} --- the association sentence with no causal
verb, and the verdict on the memo. Problem~18 is the payoff: last quarter, before the Center
opened, $45/60 = 75\%$ of the users had already failed a class against $9/90 = 10\%$ of the
non-users. The low grades sent the students to the Center; the Center did not lower the grades.

Circulate during Guided Practice and demand the reason, not the word:
\begin{itemize}[leftmargin=1.1em, nosep]
  \item \textbf{Problem 15:} ``Who put each student in their group?'' (They walked in.)
  \item \textbf{Problem 16:} ``Which whole goes with which count --- $60$ or $90$?'' A student
        who divides by $150$ has merged the two groups.
  \item \textbf{Problem 17:} circle \emph{hurts, lowers, causes} wherever they appear and hand
        the paper back. Do not explain --- just circle.
  \item \textbf{Problem 18:} ``Which came first?'' Then: ``So does the Center look better or
        worse than it really is?''
\end{itemize}
While circulating, pick the papers to put up in the debrief --- one problem~17 that wrote a
clean association sentence, and one that still has a causal verb in it. \textbf{Your formative
read comes twice}: here, and again in the last ten minutes as students start the homework.
\end{multicols}
\end{skillbox}

\boxguard[24]
\begin{skillbox}[Debrief --- whole class, spoken (10 min)]{frost}
Pens down, papers up. \textbf{This debrief is spoken} --- there is no box at the end of the
notes to fill in and no exit ticket, so it lives entirely on the board and in the room.
\begin{multicols}{2}
\begin{enumerate}[leftmargin=1.3em, nosep, itemsep=3pt]
  \item \textbf{Put the three stories back up} and make a student say the sentence: the claim,
        the arrow backwards, and retirement --- three stories, one table, and the study cannot
        tell them apart. Then state it as PS.DC.2d: an observational study shows association,
        never causation, because nobody assigned the groups.
  \item Put up the two problem~17s you picked. The difference between them is not arithmetic
        --- it is one verb. Read both aloud and ask which one the principal could act on.
  \item Take problem~18 next: which came first, and does the Center look better or worse than
        it is? This is the one that makes ``backwards'' real rather than a slogan.
  \item Close on the \textbf{banned verbs} --- \emph{causes, makes, raises, improves, helps,
        hurts} --- and problem~14: a prediction about switching groups is a causal claim.
\end{enumerate}
\columnbreak
\textbf{Check for understanding --- ask it cold, whole class.} \emph{``Bayside Middle School.
Of the $150$ students in band, $90$ earned an A or B in math; of the $600$ who are not in band,
$240$ did. The principal concludes that band \emph{raises} math grades and wants it required
for everyone. Is the principal right?''} Hands down, thirty seconds of think time, then
cold-call three students. Correct answer: $90/150 = 60\%$ against $240/600 = 40\%$, a
$20$-point gap --- \textbf{an association, and nothing more}. Nobody was placed in band;
students signed up. A concrete alternative: students who sign up for band may already study
more, or families with time for band have time for homework. \emph{This replaces the exit
ticket.}

\textbf{Formative read.} Sort the three answers as they come: \textbf{(a)} said \emph{no} and
named a concrete alternative for \emph{this} context; \textbf{(b)} said \emph{no} because ``you
can't prove causation,'' with nothing attached to the context; \textbf{(c)} said \emph{yes}, or
said the study was biased. Pile (b) is the teachable one and will be large --- it is a slogan,
not a justification, and those students abandon it the first time a causal story sounds
plausible. \textbf{If pile (c) runs past a third of the room, open Lesson~1.6 by re-running the
\emph{Not Causation} row (problems 11--14)}, and say so plainly.

\textbf{If the debrief runs short}, cut items 3 and~4 --- item~1 is the one that cannot be cut.
\textbf{Do not let the debrief eat the last ten minutes} --- the homework start is not optional
padding, it is your second formative read.
\end{multicols}
\end{skillbox}

\boxguard
\begin{skillbox}[Homework --- scored, started in class, due after two study halls]{goldbox}
Six exercises on the \textbf{Millbrook Public Library} ($1{,}200$ card holders; an unbiased
phone sample of $240$, everyone reached): \textbf{1} name the kind of study, who decided the
groups, and the explanatory and response variables; \textbf{2} the two rates with two different
denominators --- $72/90 = 80\%$ against $48/150 = 32\%$, a $48$-point gap; \textbf{3} pool the
sample and scale it up ($120/240 = 50\%$, about $600$ of $1{,}200$ --- \emph{the spiral to
Lesson~1.2}); \textbf{4} two sentences on the same data, one supported and one not --- \textbf{the
contrast pair}, the association against the newsletter's \emph{``more than doubles''};
\textbf{5} the previous summer's records ($63/90 = 70\%$ against $30/150 = 20\%$) --- the arrow
backwards, and which direction the error runs; \textbf{6} \emph{the crux head-on}: the sample
was random, everyone answered, every count is right, so explain why it \emph{still} cannot show
the program caused the extra reading. The homework is capped at \textbf{two pages}.

\textbf{Start it in the last ten minutes of the period, in class and alone.} \textbf{Scoring:}
12 points --- 2 per item. \textbf{Item~6 is where the misconception shows}: a student who
answers ``the sample was too small'' or ``it was biased'' is reaching for Lesson~1.4's tool.
\textbf{Item~4's second row carries the second check} --- the causal verb.

\textbf{Packet or Desmos --- decided here.} The packet pages are authored and are the default.
Desmos Classroom's study-design coverage is thin, and nothing there asks a student to name a
lurking variable or rewrite a causal sentence --- the items this lesson exists for --- so
\textbf{1.5 is a packet night}. Say so aloud at Close \& Assign.

\textbf{Preview of Lesson~1.6.} An observational study cannot reach causation because the
groups build themselves. So build them instead --- treatment and control, \emph{assigned} at
random --- with the four principles: comparison, randomization, replication, control.

\textbf{Connections \& Big Ideas.} \textit{Unit 1 core idea:} a conclusion is only as
trustworthy as the study that produced it. \textit{Standards covered:} PS.DC.2d (with
PS.DC.1a--b carried forward in the cycle and the conclusion items). \textit{Spirals forward
to:} experimental design (1.6, PS.DC.3a--b), comparing experiments with observational studies
(1.7, PS.DC.3c--e), the class-survey project (1.8), two-way tables (3.1), and --- directly ---
association versus causation with scatterplots and regression in Unit~5 (PS.DS.5, PS.DS.6).
\textit{Reaches back to:} Lesson~1.1's statistical cycle, Lesson~1.3's sampling plan, and
Lesson~1.4's bias, which this lesson removes on purpose so that what remains is visible.
\end{skillbox}

\boxguard
\begin{skillbox}[Watch For (while circulating)]{redbox}
\begin{itemize}[leftmargin=1.2em, nosep]
  \item \textbf{``The study was biased''} as the all-purpose answer (problems~13 and~17, HW~6).
        It is Lesson~1.4's tool applied to a lesson that deliberately removed bias. Probe:
        \emph{``The sample was random and everyone answered --- so what is left?''}
  \item \textbf{Association read as causation} (problem~13, problem~17, HW~4 --- the target
        misconception). Probe: \emph{``How many different stories fit this one table? Name the
        other two.''}
  \item \textbf{A causal verb sneaking into a written answer} (problems~14 and~17, HW~4).
        Circle \emph{causes, makes, raises, improves, helps, hurts} and hand the paper back
        without explaining. Most students fix it themselves on the second pass.
  \item \textbf{A ``lurking'' variable that is really one of the two variables} (problem~12,
        HW~5). ``Study Center users study less'' is not a third variable. Test: \emph{``Is it
        in the table? If yes, it is not lurking.''}
  \item \textbf{The reversed arrow missed entirely} (problems~11 and~18, HW~5). Most students
        accept the direction they are handed. Probe: \emph{``Which came first?''}
  \item \textbf{Two denominators merged into one} (problems~9 and~16, warm-up item~3, HW~2):
        $39\div 150$ instead of $39 \div 60$, or averaging $70$ and $45$. Probe: \emph{``Which
        whole goes with that count?''}
  \item \textbf{Sorting by ``is it a survey'' instead of ``who assigned''} (problem~3, the
        trap). Probe: \emph{``Who put that person in their group?''}
  \item \textbf{``No'' with no reason} (the cold check, HW~6). Half credit. The standard is met
        when the student names a specific alternative explanation for \emph{this} context.
\end{itemize}
Cold-call prompts: \emph{``Who decided the groups?''} \quad \emph{``Which came first?''} \quad
\emph{``What is in this study, and what is outside it?''} \quad \emph{``Say that again without
the word \emph{causes}.''}
\end{skillbox}

\boxguard
\begin{skillbox}[Close \& Assign (10 min)]{goldbox}
\textbf{Students start the homework here, in class, alone and silent --- this is the individual
practice, and the first minutes of it are your best formative read.} Hand the launch first ---
six exercises on paper, scored out of 12, due at the start of the first class after two study
halls --- and point out that the \emph{Keep in Mind} box on the cover is the page to flip back
to, because it carries the banned-verb list. Circulate: \textbf{item~6 is the column that tells
you whether the \emph{Not Causation} row landed}, and it is at the bottom of page~2, so put
students on 1, 2, and~4 first and let 3, 5, and~6 go home. Sort what you see into three piles
as you walk --- named a concrete alternative explanation (ready); said \emph{no} with only the
slogan (ready, needs the language); said the study was biased or accepted the claim (the
misconception survived) --- and open Lesson~1.6 by re-running problems 11--14 if that third
pile ran past a third of the room. Then close on what changed: \emph{``You walked in able to
turn $56$ of $80$ into a percent. You walk out knowing that a $25$-point gap between two groups
that built themselves is an \textbf{association} --- and that a random sample, a $100\%$
response rate, and perfect arithmetic still do not buy you the word \emph{causes}.''}
\textbf{Preview:} Lesson~1.6, \emph{Experimental Design} --- today nobody assigned the groups,
so we could not tell three stories apart. Next time we assign them ourselves.
\end{skillbox}

% ── Teacher notes — one per component, in packet order (three) ──────────────
% Teacher-only prose lives HERE, never in a _key: a note is the one block in a key with no
% counterpart in the blank. No note for the debrief or the homework start — both are fully
% specified in their own boxes above.
\begin{teachernote}[Warm-Up]
Five minutes, silent, timed. All three items are Algebra~1 percent work and the class has run
the identical moves in 1.2 and 1.4, so this should be near-universal --- say so, so nobody
thinks they are being retested.

\textbf{The one thing to watch is the denominators.} Item~1 divides by $80$ and item~2 divides
by $120$ --- different groups, different wholes. Item~3's third row then asks for the pooled
rate, and the wrong answer is $(70+45)\div 2 = 57.5\%$, which averages two unequal groups as if
they were the same size; the right path is $110 \div 200 = 55\%$. Catch it at the collection,
because row~3 of the notes and the entire homework depend on students keeping two denominators
straight.

\textbf{Leave $70\%$, $45\%$, $25$ points, and $55\%$ on the board and do not erase them.} The
hook needs them thirty seconds after the warm-up is collected, problem~9 reads them back off
the two-way table, and problem~10 scales the $55\%$ to $825$ of all $1{,}500$ members.
\end{teachernote}

\begin{teachernote}[Guided Notes \& Practice]
\textbf{This one document is the lesson --- pace it 20 / 15, then 10 for the debrief, then 10
for the homework start.} It is one table, and every row runs the same way: you read the
definition, mark up the display, and work the first problem of the grid (1, 5, 9, 11); the
class works the rest with the pen in their hand. Rows 1--3 must move: four, five, and four
minutes. \textbf{Take the hand count on problem~3 before anyone writes} --- it is the trap, the
one plan where the director assigns, and students who write first sort by ``is it a survey''
and never feel it.

\textbf{Row~2 is where PS.DC.2d is actually met} --- the standard asks students to \emph{use the
statistical cycle to plan and conduct an observational study}, and that row is the plan. Do not
skip it because it looks like review of 1.1. The Collect stage should come back word-for-word
from Lesson~1.3 (a random sample of $200$ from the full membership list), and it is worth
saying out loud that the plan is 1.3's and the bias check is 1.4's, so everything the class
already knows is still standing. Problem~8 is the row's punchline: the absence of any
instruction about when to swim \emph{is} the definition.

\textbf{Spend the time in row~4, \emph{Not Causation}}, which carries the misconception, and do
not reveal the retirement row too early. Give story~1, then ask the class to invent story~2
before you offer it; somebody reliably says ``maybe they get up early because they sleep
well,'' and the lesson is much stronger when it comes from the room. \emph{Only then} show
$75\%$ against $15\%$ and let the class compute it at problem~12. \textbf{Re-take the hook vote
at problem~13 before you say anything.} Let the room argue, then land the printed sentence word
by word. \textbf{Do not skip problem~14} --- the prediction about switching is the version of
this error that survives the lesson and resurfaces in every conclusion they write.

\textbf{Guided Practice (15--18) is where the pen changes hands.} \emph{Study Center} runs the
four moves the homework will ask alone --- identify, compute, write the association sentence,
test the direction --- on a school instead of a pool. It is built so that the obvious reading
is exactly backwards, and the class needs to feel the pull of the wrong answer first: let the
memo sit unchallenged until problem~17. Problem~18 is the best item on the page, and the answer
--- \emph{the Center looks worse than it is} --- requires chaining last quarter's $75\%$ against
$10\%$ to the comparison. If you only get through 15, 16, 17, and~18's verdict, that is the
right trade; what would settle it can be said aloud and is Lesson~1.6.

\textbf{There is no independent practice set on this page, and that is deliberate --- the
homework is the individual practice.} Guided Practice is the last row; the period ends with
students starting the homework in front of you, which is where your second formative read comes
from. \textbf{Do not let the debrief eat the last ten minutes.}
\end{teachernote}

\boxguard[26]
\begin{teachernote}[Homework]
\textbf{Scored, 12 points (2 per item), due at the start of the first class after two study
halls.} The ten minutes at the end of the period are a \emph{start}, not a completion: put
students on 1, 2, and~4 while you can still catch a wrong start, and let 3, 5, and~6 go home.

Items 2 and~3 are the arithmetic spine and are designed to be done in order: $80\%$, $32\%$,
the $48$-point gap, then the pooled $50\%$ and $600$. The predictable error is $120/1200$ or
averaging $80$ and $32$; both come from losing track of which whole is which --- the same error
the warm-up screens for.

\textbf{Item~4 is the fullest diagnostic in the set.} The first row is the association and is
supported; the second is the newsletter's \emph{``more than doubles how much our members
read''} and is not. Score the \emph{reason}, not the verdict --- a student who marks both
correctly and writes ``you can't prove causation'' for the second has the slogan without the
justification and will lose it on the unit test.

\textbf{Item~6 is the closer and the whole point of putting this lesson after 1.4.} The sample
was random, everyone answered, every count is right --- so the answer cannot be ``the study was
biased.'' It has to be that the card holders sorted \emph{themselves}, and no sampling method
fixes that. Write the correction on the paper rather than just marking it wrong; the same
reflex costs them on the unit test. \textbf{Sort item~6 into three piles when you score} ---
named a concrete alternative (ready); the slogan alone (ready, needs the language); blamed bias
or the sample size (the misconception survived). If more than a third land in the last pile,
reopen the three stories in the Lesson~1.6 warm-up debrief rather than letting it ride.

\textbf{Packet is the call for 1.5.} Desmos carries percent arithmetic and nothing in this set
that matters; do not swap it in. The closing \textbf{spiralbox} hands off to Lesson~1.6 and
names the four principles of experimental design --- do not teach them here.
\end{teachernote}
\end{document}
